\section{Exercício Prático: WordPress + MySQL}

\begin{frame}
  \frametitle{Exercício: WordPress + MySQL}
  \textbf{Objetivo}: Montar um ambiente completo de WordPress com banco MySQL
  usando Docker Compose, com rede para comunicação automática.

  \vspace{0.5em}
  Estrutura do projeto:

  \vspace{0.3em}
  \texttt{wordpress-docker/}

  \texttt{  compose.yaml}
\end{frame}

\begin{frame}[fragile]
  \frametitle{compose.yaml — Serviço MySQL}
  \small
\begin{verbatim}
services:
  db:
    image: mysql:8.0
    restart: always
    environment:
      MYSQL_ROOT_PASSWORD: senharoot
      MYSQL_DATABASE: wordpress
      MYSQL_USER: wpuser
      MYSQL_PASSWORD: wpsenha
    volumes:
      - db_data:/var/lib/mysql
    networks:
      - wp-network
\end{verbatim}
\end{frame}

\begin{frame}[fragile]
  \frametitle{compose.yaml — Serviço WordPress}
  \small
\begin{verbatim}
  wordpress:
    image: wordpress:latest
    restart: always
    ports:
      - "8080:80"
    environment:
      WORDPRESS_DB_HOST: db:3306
      WORDPRESS_DB_USER: wpuser
      WORDPRESS_DB_PASSWORD: wpsenha
      WORDPRESS_DB_NAME: wordpress
    volumes:
      - wp_data:/var/www/html
    networks:
      - wp-network
    depends_on:
      - db

volumes:
  db_data:
  wp_data:

networks:
  wp-network:
\end{verbatim}
\end{frame}

\begin{frame}
  \frametitle{O que está acontecendo?}
  \begin{itemize}
    \item \textbf{db}: MySQL com dados persistidos via volume \texttt{db\_data}
    \item \textbf{wordpress}: Porta 8080, conecta ao MySQL pelo hostname \texttt{db}
    \item \texttt{depends\_on}: MySQL inicia antes do WordPress
    \item \texttt{wp-network}: Rede customizada com DNS automático
    \item WordPress acessa MySQL por \texttt{db:3306} — Docker resolve o nome
  \end{itemize}
\end{frame}

\begin{frame}
  \frametitle{Executando}
  Subir toda a stack:

  \vspace{0.3em}
  \texttt{docker compose up -d}

  \vspace{0.5em}
  Verificar serviços:

  \vspace{0.3em}
  \texttt{docker compose ps}

  \vspace{0.5em}
  Acessar: \texttt{http://localhost:8080}

  \vspace{0.5em}
  \begin{block}{Resultado}
    Assistente de instalação do WordPress exibido.
    Todo o ambiente iniciado com \textbf{um único comando}.
  \end{block}
\end{frame}

\begin{frame}
  \frametitle{Derrubando o Ambiente}
  Parar e remover (volumes preservados):

  \vspace{0.3em}
  \texttt{docker compose down}

  \vspace{0.8em}
  Parar, remover E apagar volumes:

  \vspace{0.3em}
  \texttt{docker compose down -v}

  \vspace{0.8em}
  \begin{alertblock}{Cuidado}
    A flag \texttt{-v} remove os volumes, apagando \textbf{todos os dados}.
    Sem ela, os dados estarão lá na próxima vez.
  \end{alertblock}
\end{frame}
